% GENERATED FILE. Edit canonical lessons, then run npm run generate:books.
% canonical-source-hash: fnv1a64:b767ee4305dec7c3
% canonical-lessons: MR-C30-aani, MR-C30-kinva, MR-C30-pan, MR-R30-three-joins

\chapter{Two Sentences, One Breath}
\label{ch:mr-flat-joins}
\hlchaptermodality{30}

\begin{chapteropening}
\textbf{By the end of this chapter:} \emph{I can join two Marathi sentences with aani, offer a choice with kinvaa, and take something back with pan.}

Complete MR-R30-three-joins and fill the same slot three ways with chaha and paani: added, offered, contradicted.
\end{chapteropening}

\section[āṇi]{\mr{आणि} (āṇi) — and}

\label{lesson:MR-C30-aani}



\begin{quote}
Say two sentences you already own: \textbf{\mr{मला} \mr{चहा} \mr{आवडतो}}, and \textbf{\mr{मला} \mr{दूध} \mr{आवडतं}}. Two sentences, two breaths. This chapter makes them one.
\end{quote}

\subsection*{You'll want to know: \mr{आणि}}
\begin{quote}
\textbf{\mr{आणि}} — \emph{āṇi} — \textbf{and}
\end{quote}

\begin{quote}
\textbf{\mr{मला} \mr{चहा} \mr{आवडतो} \mr{आणि} \mr{दूध} \mr{आवडतं}.} — \emph{malā chahā āvaḍto āṇi dūdh āvaḍtaṁ.} — \textbf{I like tea and I like milk.}
\end{quote}

\textbf{\mr{आणि}} sits BETWEEN the two pieces and changes nothing on either side. Each half keeps its own agreement: \textbf{\mr{चहा}} is masculine and keeps \textbf{\mr{आवडतो}}, \textbf{\mr{दूध}} is neuter and keeps \textbf{\mr{आवडतं}}. The joining word does not reach inside.

It joins two nouns just as happily: \textbf{\mr{भाऊ} \mr{आणि} \mr{बहीण}}, brother and sister.

The word is old and slightly surprising. It comes from Sanskrit \textbf{\mr{अन्य}} (\emph{anya}), \textquotedblleft{}\textbf{other}.\textquotedblright{} A word for \emph{the other one} wore down into a word for \emph{one more}. English does the same trick from the other side: \emph{another} is \emph{an other}.

\subsection*{Your turn}
Say these aloud:

\begin{itemize}
\raggedright
  \item \emph{bhāū āṇi bahīṇ} — brother and sister
  \item \emph{malā chahā āvaḍto āṇi dūdh āvaḍtaṁ}
  \item which half of that sentence changed when you joined it — \textbf{neither}
\end{itemize}

\begin{tcolorbox}[breakable,colback=teal!4,colframe=teal!35!black,title={Before you move on}]
Where does \textbf{\mr{आणि}} sit, and what does it change? (\textbf{Between the two pieces; nothing}.) What did it originally mean? (\textbf{Other}.)

Sources: \href{https://en.wiktionary.org/wiki/\%E0\%A4\%86\%E0\%A4\%A3\%E0\%A4\%BF}{Wiktionary: \mr{आणि}}.
\end{tcolorbox}
\section[kiṁvā]{\mr{किंवा} (kiṁvā) — or}

\label{lesson:MR-C30-kinva}



\begin{quote}
One lesson ago you got \textbf{\mr{आणि}}, which adds. Say \textbf{\mr{चहा} \mr{आणि} \mr{पाणी}}. Now you need the word that does the opposite: it puts two things side by side and asks you to take only one.
\end{quote}

\subsection*{You'll want to know: \mr{किंवा}}
\begin{quote}
\textbf{\mr{किंवा}} — \emph{kiṁvā} — \textbf{or}
\end{quote}

\begin{quote}
\textbf{\mr{चहा} \mr{किंवा} \mr{पाणी}?} — \emph{chahā kiṁvā pāṇī?} — \textbf{Tea or water?}
\end{quote}

It sits exactly where \textbf{\mr{आणि}} sat. The two words are a matched pair: same slot, opposite job. \textbf{\mr{आणि}} hands you both, \textbf{\mr{किंवा}} makes you choose.

Taken apart, \textbf{\mr{किंवा}} is two Sanskrit words carried over whole: \textbf{\mr{किम्}} (\emph{kim}), \textquotedblleft{}what\textquotedblright{} — the same stem that gave Marathi \textbf{\mr{काय}} — plus \textbf{\mr{वा}} (\emph{vā}), an ancient particle meaning \textquotedblleft{}or.\textquotedblright{} That \textbf{\mr{वा}} is genuinely old: Latin kept it as the \textbf{-ve} in \emph{sīve}, \textquotedblleft{}or if.\textquotedblright{} Marathi's word for choosing is, underneath, \textquotedblleft{}\textbf{what-or}.\textquotedblright{}

\subsection*{Your turn}
Say these aloud:

\begin{itemize}
\raggedright
  \item \emph{chahā kiṁvā pāṇī?}
  \item \emph{dūdh kiṁvā pāṇī?}
  \item the same pair with \emph{āṇi}, and hear the offer become a list
\end{itemize}

\begin{tcolorbox}[breakable,colback=teal!4,colframe=teal!35!black,title={Before you move on}]
Which of the two joining words makes the listener choose? (\textbf{\mr{किंवा}}.) What two Sanskrit words is it made of? (\textbf{\mr{किम्} + \mr{वा}}, \textquotedblleft{}what or\textquotedblright{}.)

Sources: \href{https://en.wiktionary.org/wiki/\%E0\%A4\%95\%E0\%A4\%BF\%E0\%A4\%82\%E0\%A4\%B5\%E0\%A4\%BE}{Wiktionary: \mr{किंवा}}.
\end{tcolorbox}
\section[paṇ]{\mr{पण} (paṇ) — but}

\label{lesson:MR-C30-pan}



\begin{quote}
Two joining words are yours: \textbf{\mr{आणि}} adds, \textbf{\mr{किंवा}} offers. Say \textbf{\mr{चहा} \mr{आणि} \mr{दूध}}, then \textbf{\mr{चहा} \mr{किंवा} \mr{दूध}}. Both are friendly. Neither can disagree with anything.
\end{quote}

\subsection*{You'll want to know: \mr{पण}}
\begin{quote}
\textbf{\mr{पण}} — \emph{paṇ} — \textbf{but}
\end{quote}

\begin{quote}
\textbf{\mr{चहा} \mr{आहे} \mr{पण} \mr{दूध} \mr{नाही}.} — \emph{chahā āhe paṇ dūdh nāhī.} — \textbf{There is tea but there is no milk.}
\end{quote}

\textbf{\mr{पण}} stands in the same slot as the other two and does the one thing they cannot: it lets the second half push back against the first. Until you have it, every sentence you can build only agrees with itself.

\textbf{\mr{पण}} and \textbf{\mr{पुन्हा}} — the \emph{again} in \textbf{\mr{पुन्हा} \mr{भेटू}} — are the same Sanskrit word, \textbf{\mr{पुनर्}} (\emph{punar}), \textquotedblleft{}again, back.\textquotedblright{} One descendant kept the time meaning; the other kept only the \emph{back}, and turned it into contradiction. English does this too: \emph{but} began as \emph{be-utan}, \textquotedblleft{}outside.\textquotedblright{}

\subsection*{Your turn}
Say these aloud:

\begin{itemize}
\raggedright
  \item \emph{chahā āhe paṇ dūdh nāhī}
  \item the three joining words in a row — \emph{āṇi}, \emph{kiṁvā}, \emph{paṇ}
  \item which Marathi word for \emph{again} is \textbf{\mr{पण}}'s twin — \emph{punhā}
\end{itemize}

\begin{tcolorbox}[breakable,colback=teal!4,colframe=teal!35!black,title={Before you move on}]
Which joining word lets the second half disagree with the first? (\textbf{\mr{पण}}.) Which everyday farewell word is its twin? (\textbf{\mr{पुन्हा}}.)

Sources: \href{https://en.wiktionary.org/wiki/\%E0\%A4\%AA\%E0\%A4\%A3}{Wiktionary: \mr{पण}}.
\end{tcolorbox}
\section[… āṇi … · … kiṁvā … · … paṇ …]{The same gap, three ways}

\label{lesson:MR-R30-three-joins}



\begin{quote}
Three words landed in this chapter and all three go in the same place. Name them in the order you met them.
\end{quote}

\subsection*{Your turn}
Say these aloud:

\begin{itemize}
\raggedright
  \item \emph{malā chahā āvaḍto āṇi dūdh āvaḍtaṁ}
  \item \emph{chahā kiṁvā pāṇī?}
  \item \emph{chahā āhe paṇ dūdh nāhī}
\end{itemize}

\begin{tcolorbox}[breakable,colback=teal!4,colframe=teal!35!black,title={Before you move on}]
Which one hands over both? (\textbf{\mr{आणि}}.) Which one makes you pick? (\textbf{\mr{किंवा}}.) Which one takes something back? (\textbf{\mr{पण}}.) Then stop.
\end{tcolorbox}
